% Inhoud van één labo. Gebruikt door de syllabus van dit labo en door de syllabus voor labo 3-5.
\textbf{Doel:} een grafiektype kiezen voor een gegeven vraag, de
standaardgrafieken van \code{plotly.express} maken, extra dimensies
toevoegen met kleur, facetten en grootte, en eenvoudige interactiviteit
inbouwen.

\begin{itemize}
  \item \textbf{Principes} -- De juiste grafiek voor de vraag (verdeling, vergelijking, verband, evolutie); kritiek op een misleidende grafiek.
  \item \textbf{Basis van \code{plotly.express}} -- Histogram, boxplot, spreidingsdiagram en lijngrafiek, met titels, aslabels en hover-informatie.
  \item \textbf{Meer dimensies} -- \code{color}, \code{facet\_col} en \code{size}; één vaste kleur per positie.
  \item \textbf{Hartslag als tijdreeks} -- Zonebanden met \code{add\_hrect}, en waarom we geen tweede y-as gebruiken.
  \item \textbf{Interactiviteit} -- \code{mo.ui.dropdown} en \code{mo.ui.range\_slider} als keuzemenu en schuifbalk; optioneel selecties in een grafiek met \code{mo.ui.plotly}.
  \item \textbf{Minidashboard} -- 3 grafieken en 1 keuzemenu waarop minstens één grafiek reageert.
\end{itemize}

\textbf{Oefeningen:} (1) een histogram van het gewicht; (2) Yo-Yo-niveau
per positie, met facetten per geslacht; (3) snelheid doorheen de tijd
voor één atleet; \textbf{opdracht:} het minidashboard.
